\section{The new weighted certificate for choosing a contraction plan}
\label{sec:portfolio}

\subsection{Every matrix operation has nine linear supports}

An operation merges disjoint sets $A,B$ with outside set $O$, or contracts
the last three blocks. Its three index classes are the unordered pair sets
$AB,AO,BO$. Their logarithmic dimensions are nonnegative numbers
$H_1,H_2,H_3$. Formula~\eqref{eq:rectangle}, after sorting these numbers,
is exactly the maximum of nine linear forms on the unsorted triple:
\[
 \operatorname{Perm}(1+\varepsilon/2,1+\varepsilon/2,0)
 \quad\text{and}\quad
 \operatorname{Perm}(1,47/56+\varepsilon,15/28).
\]
The first set has three vectors and the second six. This follows from the
rearrangement inequality, since $1\ge47/56+\varepsilon\ge15/28$.
They represent one sorted cost function, rather than nine separately
executed algorithms.

Multiply the supports by the fixed denominator $L_*=1{,}400{,}000$.
The integer triples become
\begin{equation}\label{eq:integer-triples}
 (1400007,1400007,0),\qquad(1400000,1175014,750000),
\end{equation}
with the permutations just specified. Assign each unordered group edge its
class coefficient, and assign zero to edges internal to a part.

A plan is an ordered partition of $[r]$ into three nonempty blocks. Build
each block by splitting its ordered list at $\lfloor k/2\rfloor$,
recursively building the two halves, then multiplying their tables.
Finally contract the three block tables. There are $r-3$ internal merges
and one final triangle, hence $r-2$ operations and $9(r-2)$ supports.
For plan $j$ let $\mathcal V_j$ be this finite set, and put
\[
 F_j(g)=L_*^{-1}\max_{v\in\mathcal V_j}v\cdot g.
\]
Up to polynomial factors, $2^{F_j(g)}$ pays for every local operation
after the singleton group tables have been computed.

\subsection{Couple the norm terms before taking pairwise maxima}

The preceding tree packet used all tuples of one support from each plan.
With $q$ plans and at most $m$ supports per plan, that exact certificate
has $m^q$ tuples. The following certificate needs only $\binom q2m^2$
support pairs. Its inequality can be weaker than complete tuple
enumeration, but its manageable cost permits a larger, more effective
menu of plans.

\begin{theorem}[Weighted pair certificate]\label{thm:pair}
Let $q\ge2$ and fix nonnegative integer weights $w_1,\ldots,w_q$ with
$W=\sum_iw_i>0$. For $i<j$ define
\begin{align}
 A_{ij}=\max_{\substack{v\in\mathcal V_i\\z\in\mathcal V_j}}
 \bigl(&w_i^2\norm v^2+w_j^2\norm z^2\notag\\
       &+2(q-1)w_iw_j\,v\cdot z\bigr),\label{eq:pair-max}
\end{align}
and
\[
 U=\frac{\sum_{i<j}A_{ij}}{(q-1)W^2},\qquad
 \beta^2=\frac{2U}{L_*^2r^2}.
\]
Selecting the plan of least $F_j(g)$ gives local nonrecursive time and
local storage at most a polynomial factor times
$2^{(\beta+o(1))\K(M,u)}$, uniformly over all exact mode-budget
allocations.
\end{theorem}

\begin{proof}
For arbitrary choices $v_i\in\mathcal V_i$, expansion of the squared norm
gives the exact identity
\begin{align*}
 \left\|\sum_i\frac{w_i}{W}v_i\right\|_2^2
 =\frac1{(q-1)W^2}\sum_{i<j}
 \left(w_i^2\norm{v_i}^2+w_j^2\norm{v_j}^2
       +2(q-1)w_iw_jv_i\cdot v_j\right).
\end{align*}
Each diagonal term occurs $q-1$ times. Taking the maximum of each pair
bracket separately can only increase the sum, so the squared norm is at most $U$.
For a given $g$, choose a maximizing support $v_i$ for each $F_i(g)$.
Then
\[
 \min_iF_i(g)\le\sum_i\frac{w_i}{W}F_i(g)
 =\frac1{L_*}\left(\sum_i\frac{w_i}{W}v_i\right)\cdot g
 \le(\beta+o(1))\K
\]
by~\eqref{eq:undirected-entropy}. The fixed finite menu makes the error
uniform. There is no exchange of a minimum and a maximum, and no assumption
that independently maximizing supports can be realized simultaneously.

Every table is an input or output of one of the plan's paid matrix
operations. Its logarithmic size is a sum of two dimension logarithms,
which is bounded by the first branch of~\eqref{eq:rectangle}. The
backend's space bound therefore covers all local tables, including the
singleton caches. Their costly recursive evaluation time is separate and
remains in~\eqref{eq:recurrence}.
\end{proof}

The weights are a proof device, not random choices at runtime. The
algorithm always selects the cheapest actual plan. Compared with separately
bounding every norm and inner product, \eqref{eq:pair-max} also keeps the
correlation between the norm and dot product of each pair of supports.

\subsection{The complete ten-group certificate}

Group labels start at zero. Each ordered block below is split as above.
The plan weights have denominator $W=1000$.

\begin{center}\small
\begin{tabular}{crlll}
\toprule Plan & $w_j$ & First block & Second block & Third block\\\midrule
1&137&$(7,4,3,9)$&$(1,6,0)$&$(8,2,5)$\\
2&124&$(5,4,6,7)$&$(9,1,8)$&$(2,3,0)$\\
3&135&$(0,5,9,7)$&$(2,1,4)$&$(8,3,6)$\\
4&120&$(0,4,7,8)$&$(2,9,6)$&$(1,5,3)$\\
5&108&$(3,1,7,2)$&$(9,6,8)$&$(5,4,0)$\\
6&119&$(7,1,5,8)$&$(2,9,4)$&$(0,6,3)$\\
7&106&$(7,2,8,6)$&$(5,9,3)$&$(0,1,4)$\\
8&87&$(1,9,6,7)$&$(2,5,0)$&$(4,3,8)$\\
9&64&$(4,6,1,5)$&$(8,9,2)$&$(0,7,3)$\\
\bottomrule
\end{tabular}\end{center}

There are 72 supports per plan, each in dimension $\binom{10}2=45$.
The exact checker reconstructs all 648 supports and evaluates
$\binom92\,72^2=186{,}624$ pair choices. It obtains
\begin{equation}\label{eq:beta10}
 \boxed{\displaystyle\beta_{10}^2=\CachedBetaSquared.}
\end{equation}
The JSON result contains every exact pair maximum and a maximizing support
pair, so the certificate does not depend on a floating-point search result.
The finite sum is reproduced from the plans and integer triples alone.

\subsection{Exact selection and complexity of the selector}

For each plan compute the integer
\[
 P_j=\max_{v\in\mathcal V_j}
 \prod_{\{g,h\}}
       \left(|\mathcal I_{gh}|\,|\mathcal I_{hg}|\right)^{v_{\{g,h\}}}.
\]
Minimizing $P_j$ is exactly equivalent to minimizing $F_j(g)$. Its bit
length is polynomial in $s$: $r$, $q$, the support coefficients and their
denominator are fixed constants, and the cardinalities have polynomial bit
length. Thus plan selection has no logarithmic comparison oracle or
rounding ambiguity. The literal integer-product routine in the code is
intended as an exact specification; the large fixed denominator makes it
unsuitable as a claim of practical efficiency.

The improvement is new relative to the reviewed corpus: a fixed finite
matrix-operation portfolio is certified by weighted coupled pairs, and
that certificate is integrated with a paid recursive schedule. Averaging,
matrix multiplication, and quadratic norm identities are classical; no
literature-wide priority claim is made for those ingredients.
